\section{Log-Gamma integrals, logarithmic derivatives, and the improper integral}\label{sec:integrals}

\subsection{Alexeiewsky's log-Gamma integral}

The DLMF form most convenient for branches is
\begin{equation}\label{eq:alex-dlmf}
 \Log G(1+z)=\frac z2\log(2\pi)-\frac{z(z+1)}2
 +z\Log\Gamma(1+z)-\int_0^z\Log\Gamma(1+t)\dd t.
\end{equation}
Here the logarithms are continued from the positive real axis along a path
avoiding the poles and zeros.  Replacing \(t+1\) by \(x\), or using the Gamma
functional equation, gives the equivalent target-page form
\begin{equation}\label{eq:alex}
 \int_0^z\Log\Gamma(x)\dd x
 =\frac{z(1-z)}2+\frac z2\log(2\pi)
 +(z-1)\Log\Gamma(z)-\Log G(z).
\end{equation}

\begin{theorem}[Constructive Alexeiewsky identity]\label{thm:alex}
On any rational polygonal path from a positive base point to \(z\), separated
from the nonpositive integers and equipped with coherent branches of
\(\Log\Gamma\) and \(\Log G\), \cref{eq:alex-dlmf,eq:alex} hold with an
explicit finite quadrature modulus.  Their exponentiated derivatives extend to
the corresponding meromorphic identities globally.
\end{theorem}

\begin{proof}
Define the difference between the two sides of \cref{eq:alex-dlmf} as
\(A(z)\).  Differentiate on a zero-free disk.  The canonical-product series
may be differentiated by \cref{eq:qprime-bound}.  The reciprocal-Gamma product
may be differentiated using its analogous \(k^{-2}\) tail from the Gamma
packet.  Termwise cancellation gives \(A'(z)=0\).  At \(z=0\), every term is
zero, so \(A(0)=0\).  Subdivide the certified path into disks on which both
Taylor contracts are valid.  On each segment the finite integral of the
derivative is enclosed by a rational ball of radius smaller than the allocated
budget.  A finite telescope proves \(A(z)=0\) at the requested precision.

For a more algebraic verification, start with the finite product logarithms.
After integrating the finite reciprocal-Gamma logarithm, every summand is an
explicit polynomial-log expression.  Subtract the finite Barnes logarithm;
all terms cancel except an endpoint error.  The Gamma and Barnes tail bounds
make that error tend effectively to zero.  This second route is especially
suitable for an object-language proof because it avoids a general fundamental
theorem of calculus: only a finite antiderivative identity and two tail
estimates are used.
\end{proof}

On the positive real axis, the logarithms are ordinary real logarithms and no
branch data are needed.  The integral is implemented by the fixed Gamma
quadrature contract from the existing packet.  For complex paths, the integral
is a sum over finitely many line segments and the same evaluator handles the
real and imaginary parts after an affine substitution.

\subsection{Logarithmic derivative and derivative formulas}

Differentiate \cref{eq:alex}.  The \(\log\Gamma(z)\) terms cancel and yield
\begin{equation}\label{eq:phi-global}
 \frac{G'(z)}{G(z)}
 =(z-1)\psi(z)-z+\frac12\log(2\pi)+\frac12,
\end{equation}
where \(\psi=\Gamma'/\Gamma\).  Multiplying by \(G(z)\) gives the globally
entire identity
\begin{equation}\label{eq:Gprime}
 G'(z)=G(z)\left((z-1)\psi(z)-z+\frac12\log(2\pi)+\frac12\right),
\end{equation}
understood at Gamma's poles by cancellation.  Shifting \(z\mapsto z+1\) and
using \(\psi(z+1)=\psi(z)+1/z\) gives the MathWorld form
\begin{equation}\label{eq:phi-shift}
 \frac{G'(z+1)}{G(z+1)}
 =\frac12\log(2\pi)+\frac12-z+z\frac{\Gamma'(z)}{\Gamma(z)}.
\end{equation}

\begin{proof}[Division-free formal form]
On a compact domain where \(G\) and \(\Gamma\) are apart from zero, the quotient
identities follow from the named reciprocal evaluators.  The primary theorem,
however, is obtained by multiplying \cref{eq:Gprime} by \(\Gamma(z)\):
\[
 \Gamma(z)G'(z)
 =G(z)\left((z-1)\Gamma'(z)
 +\left[-z+\tfrac12\log(2\pi)+\tfrac12\right]\Gamma(z)\right).
\]
Every factor has a Taylor-ball name, and no division appears.  Around a
nonpositive integer, use the known zero orders of \(G\) and \(1/\Gamma\) to
verify the removable singularity by finite polynomial factorization.  This is
the object-iPRA target.
\end{proof}

The Taylor series \cref{eq:phi-taylor} is recovered by differentiating the
canonical series, so the integral, product and power-series derivations are
mutually consistent without circularity.

\subsection{The improper-integral representation}

The target page gives
\begin{equation}\label{eq:improper}
 \log G(1+z)
 =\frac z2\log(2\pi)
 +\int_0^\infty\frac{1}{t}
 \left[
 \frac{1-\e^{-zt}}{4\sinh^2(t/2)}
 +\frac{z^2}{2}\e^{-t}-\frac zt
 \right]\dd t.
\end{equation}
The literal integral converges for
\begin{equation}\label{eq:improper-domain}
 \Re z>-1.
\end{equation}
At zero the three apparent singular terms cancel; at infinity the first two
terms decay exponentially and the last is \(O(t^{-2})\) after the exterior
factor \(1/t\).  If \(\Re z\le-1\), the expression may still represent the
analytic continuation, but not as this unmodified convergent improper integral.

\begin{theorem}[Improper-integral formula]\label{thm:improper}
On every compact set
\[
 |z|\le R,\qquad \Re z\ge-1+\delta,
\]
with positive rational \(\delta\), \cref{eq:improper} has a
primitive-recursive quadrature and tail modulus and equals the canonical
product definition of \(G\).
\end{theorem}

\begin{proof}
Use the exact expansion
\begin{equation}\label{eq:sinh-kernel}
 \frac{1}{4\sinh^2(t/2)}
 =\frac{\e^{-t}}{(1-\e^{-t})^2}
 =\sum_{k=1}^{\infty}k\e^{-kt},\qquad t>0.
\end{equation}
To avoid separately divergent Frullani integrals, differentiate the right side
of \cref{eq:improper} three times with respect to \(z\).  The polynomial
regularizers disappear and dominated differentiation gives
\begin{align}
 I'''(z)
 &=\int_0^\infty
   \frac{t^2\e^{-zt}}{4\sinh^2(t/2)}\dd t\notag\\
 &=\sum_{k=1}^{\infty}k\int_0^\infty t^2\e^{-(k+z)t}\dd t
 =2\sum_{k=1}^{\infty}\frac{k}{(k+z)^3}.
 \label{eq:improper-third}
\end{align}
The same third derivative follows from the product logarithm by
\cref{eq:qhigher} with \(r=3\).  The value and the first two derivatives at
\(z=0\) agree by direct finite evaluation of the regularized integral and the
coefficients in \cref{eq:taylor-logG}.  Therefore the two analytic names agree
on the connected half-plane.

For the quantitative version, split the integral into
\([0,2^{-M}]\), \([2^{-M},T]\), and \([T,\infty)\).  At zero, expand
\(1-\e^{-zt}\), \(\e^{-t}\), and
\(t^2/(4\sinh^2(t/2))\) to a fixed order.  Their constant and linear
singularities cancel exactly; the remainder is bounded by a rational multiple
of \(t\), uniformly for \(|z|\le R\).  Thus the first piece is at most
\(C_R2^{-M}\).

For the tail, \cref{eq:sinh-kernel} and
\(\Re z\ge-1+\delta\) give
\[
 \left|\frac{1-\e^{-zt}}{4\sinh^2(t/2)}\right|
 \le C_R\bigl(\e^{-t}+\e^{-\delta t}\bigr),
 \qquad t\ge1.
\]
Together with the other terms, the tail is bounded by
\[
 C_R\left(\frac{\e^{-\delta T}}{\delta T}
          +\frac{\e^{-T}}{T}+\frac1T\right).
\]
Choose primitive-recursive \(M,T\) to make the endpoint pieces smaller than
the allocated dyadic budgets.  The middle integral is ordinary rational
quadrature with a derivative modulus.  These estimates also justify the
interchanges in \cref{eq:improper-third}.
\end{proof}

\subsection{A practical evaluator}

For small \(|z|\), the Taylor series is fastest.  On moderate compact disks the
canonical product is simplest.  In the half-plane \(\Re z>-1+\delta\), the
integral can be useful when its smooth middle interval is short.  For large
\(|z|\), the asymptotic module in \cref{sec:asymptotic} dominates.  A formal
evaluator need not choose among them by an undecidable real comparison: it can
use rational upper bounds on the input box and a fixed priority table.  The
certificate records which algorithm was chosen, and the soundness theorem for
that branch supplies the result.

\begin{roadmap}[Lean modules \code{BarnesG/Integral} and
\code{BarnesG/Derivative}]
Reuse the Gamma packet's path integration and digamma contracts.  Prove the
finite version of \cref{eq:sinh-kernel} with a geometric tail.  Encode the
zero-end regularizer as a fixed Taylor polynomial plus a rational remainder.
The derivative theorem should first be cross-multiplied by \(\Gamma G\); add
quotient corollaries only under apartness structures.
\end{roadmap}
